\section{本讲小结}
牛顿运动定律把同一研究对象的外力同其加速度联系起来。先确定惯性参考系和系统边界，才能判断哪些力进入方程；再把接触、绳长和共同运动条件写成加速度关系。整体法消去的是内部成对的作用力，不是把内部物体的不同加速度合并为同一个数。

支持力、张力和静摩擦都是随运动条件确定的量。求出它们后，应检查支持力与张力的非负性，以及静摩擦上限。剪断问题则先保留连续的形变与速度，再重新求解约束力。一个合格的结果还应说清方向、适用区间和临界状态之后需要改变的方程。

\section{练习}
\begin{question}\begin{drill}[竖直加速中的叠块]\label{n1d1}
电梯以 $a_0$ 向上加速，地板光滑，下层板质量 $M$、上层块质量 $m$，层间最大静摩擦因数 $\mu_s$。用水平力 $F$ 拉下层板，求两者共同运动的最大 $F$ 以及地板支持力。
\fig{n1d1q}{竖直共同加速度改变层间法向力}{\begin{tikzpicture}
\draw[world](-2,0)rectangle(2,2.5);\node[draw,minimum width=2.2cm,minimum height=.6cm,anchor=south]at(0,0){$M$};\node[draw,anchor=south]at(0,.6){$m$};\draw[force](1.1,.3)--(2.6,.3)node[right]{$F$};\draw[flow](-2.5,.8)--(-2.5,2)node[above]{$a_0$};
\end{tikzpicture}}\answerlink{n1d1}
\end{drill}\end{question}

\begin{question}\begin{drill}[由秤读数反推质量差]\label{n1d2}
两物体由跨过轻定滑轮的绳连接，滑轮用弹簧秤悬挂。已知总质量 $M$、运动时秤读数 $S$，求两物体的质量及加速度大小，忽略滑轮自重，并说明实测值必须满足的条件。
\fig{n1d2q}{质量未知，读数与总质量已知}{\begin{tikzpicture}
\draw[spring](0,2.8)--(0,1.8);\draw[world](0,1.8)--(0,1.4);\draw[world](0,1.4)circle(.3);\draw[world](-.3,0)--(-.3,1.4)arc(180:0:.3)--(.3,.5);\node[draw]at(-.3,-.2){$m_1$};\node[draw]at(.3,.3){$m_2$};\node[right]at(.3,2.3){$S$};\node at(2.2,1){$m_1+m_2=M$};
\end{tikzpicture}}\answerlink{n1d2}
\end{drill}\end{question}

\begin{question}\begin{drill}[非零初速下的斜绳剪断]\label{n1d3}
定长轻绳长 $l$，连接固定悬点和质量 $m$ 的小球。考察绳与竖直夹角 $\theta$ 的时刻，球速率为 $v$，此时不再有其他约束。求径向、切向加速度及绳张力，说明为何不能直接套用静止释放的结果。
\fig{n1d3q}{径向指向悬点，速度沿圆弧切线}{\begin{tikzpicture}
\coordinate(O)at(0,2.5);\coordinate(B)at(1.2,.4);\draw[world](O)--(B)node[midway,right]{$l$};\draw[guide](O)--(0,0);\fill(B)circle(3pt);\node[above]at(O){$O$};\draw[flow](B)--(2.5,1.14)node[right]{$v$};\node at(.3,1.6){$\theta$};
\end{tikzpicture}}\answerlink{n1d3}
\end{drill}\end{question}

\begin{question}\begin{drill}[规定加速度的光滑斜面]\label{n1d4}
斜面向右上方升高，倾角 $\alpha$，在外力作用下以规定加速度 $a_0$ 向左平移。质量 $m$ 的物块与斜面光滑接触，求支持力及相对斜面的沿面加速度，讨论物块可能相对静止的条件。
\fig{n1d4q}{已知斜面平移，不等于已知物块完整加速度}{\begin{tikzpicture}
\draw[world](-2,0)--(2,0)--(2,2.2)--cycle;\node[draw,rotate=28.8108,anchor=south]at(.4,1.32){$m$};\draw[flow](-1.5,.6)--(-3,.6)node[left]{$a_0$};\node at(-1,.2){$\alpha$};
\end{tikzpicture}}\answerlink{n1d4}
\end{drill}\end{question}
